\documentclass{article}

\usepackage{amsmath,amssymb}
\usepackage{xurl}
\usepackage[hidelinks]{hyperref}

% Shared commands for notes in docs/paper-gaps/.

\DeclareUnicodeCharacter{2080}{\ifmmode{}_{0}\else\textsubscript{0}\fi}
\DeclareUnicodeCharacter{2081}{\ifmmode{}_{1}\else\textsubscript{1}\fi}
\DeclareUnicodeCharacter{2082}{\ifmmode{}_{2}\else\textsubscript{2}\fi}
\DeclareUnicodeCharacter{2099}{\ifmmode{}_{n}\else\textsubscript{n}\fi}

\newcommand{\Lean}{\textsc{Lean}}
\ExplSyntaxOn
\NewDocumentCommand{\leanid}{m}{
  \ifmmode
    \textsf{\detokenize{#1}}
  \else
    \group_begin:
    \urlstyle{sf}
    \tl_set:Nn \l_tmpa_tl {#1}
    \tl_replace_all:Nnn \l_tmpa_tl {\_} {_}
    \exp_args:NV \path \l_tmpa_tl
    \group_end:
  \fi
}
\ExplSyntaxOff
\providecommand{\tr}{\operatorname{tr}}
\newcommand{\tnleanIssueBase}{https://github.com/LionSR/TNLean/issues/}
\newcommand{\tnleanPrBase}{https://github.com/LionSR/TNLean/pull/}
\newcommand{\tnleanDocBase}{https://sirui-lu.com/TNLean/docs/}
\newcommand{\ghissue}[1]{\href{\tnleanIssueBase#1}{GitHub issue \##1}}
\newcommand{\ghpr}[1]{\href{\tnleanPrBase#1}{PR \##1}}
\newcommand{\doclink}[1]{\href{\tnleanDocBase#1.html}{\path{#1}}}

% Dirac notation and the identity operator, as in blueprint/src/macros/common.tex.
% \providecommand keeps note-local definitions authoritative.
\usepackage{dsfont}
\providecommand{\ket}[1]{|#1\rangle}
\providecommand{\bra}[1]{\langle #1|}
\providecommand{\braket}[2]{\langle #1 | #2 \rangle}
\providecommand{\ketbra}[2]{|#1\rangle\!\langle #2|}
\providecommand{\Id}{\mathds{1}}
\providecommand{\one}{\mathds{1}}
\providecommand{\C}{\mathbb{C}}
\providecommand{\MN}[1]{M_{#1}(\C)}
\providecommand{\MD}{M_D(\C)}

% Verdict marker: \gapnote{<kind>}{<status>}. Kind is the policy
% classification (clarification, local-correction, scope-restriction,
% unfaithful, false-source, open-gap); status is open, wip (elimination
% underway), resolved, or historical. Machine-read by the site index;
% prints a verdict line.
\newcommand{\gapnote}[2]{\par\noindent\textbf{Verdict:} #1 (#2).\par}


\title{Boundary-Dressed Gauge Scalars and L-Symbols}
\author{}
\date{2026-10-30}

\begin{document}
\maketitle
\gapnote{local-correction}{resolved}

\section{The printed statements}

Let $g\mapsto O_g$ be a group of matrix product operators and let $A_x$,
$x\in\mathsf X$, be normal matrix product state tensors permuted by it.
Fusion tensors $F^<_{g,h},F^>_{g,h}$ reduce the stacked product $O_gO_h$ onto
$O_{gh}$, and action tensors $A^\Gamma_{g,x},A^{\text{ref}\Gamma}_{g,x}$
reduce $O_g\cdot A_x$ onto $A_{gx}$.  Ref.~\cite{FrancoRubio2025MPUGauging}
states that these tensors are unique up to scalars,
\footnote{Source coordinates in \path{Papers/2502.20257/main.tex}:
    \texttt{eq:scalar\_fus\_ten} at lines~1500--1504 and
    \texttt{eq:scalar\_act\_ten} at lines~1871--1873.}
\begin{align}
    F^>_{g,h}
        &\mapsto\beta_{g,h}F^>_{g,h},
    &F^<_{g,h}
        &\mapsto\beta_{g,h}^{-1}F^<_{g,h},
    \label{eq:fbc25_dressed_fusion_gauge}\\
    A^{\text{ref}\Gamma}_{g,x}
        &\mapsto\gamma_{g,x}A^{\text{ref}\Gamma}_{g,x},
    &A^{\Gamma}_{g,x}
        &\mapsto\gamma_{g,x}^{-1}A^{\Gamma}_{g,x},
    \label{eq:fbc25_dressed_action_gauge}
\end{align}
and defines the $L$-symbol $L^x_{g,h}\in\mathbb C^\times$ as the scalar
relating the two contractions of fusion and action tensors that reduce
$(O_gO_h)\cdot A_x$ onto $A_{(gh)x}$.
\footnote{\texttt{eq:defL} at lines~1875--1913.}
It then asserts the gauge law
\begin{align}
    L^x_{g,h}
        &\mapsto\frac{\gamma_{gh,x}\beta_{g,h}}{\gamma_{g,hx}\gamma_{h,x}}
          L^x_{g,h}
    \label{eq:fbc25_dressed_l_gauge}
\end{align}
at line~1918 and the compatibility with the anomaly three-cocycle at
line~1923.

\section{What the reduction theorem provides}

The only uniqueness statement for reductions onto a normal tensor is
Theorem~22 of Ref.~\cite{Molnar2018SemiInjective}.
\footnote{\path{cornerproblem.tex} of
    \href{https://arxiv.org/abs/1706.07329v2}{arXiv:1706.07329v2},
    lines~3156--3162.}
If $(V,W)$ and $(V',W')$ reduce $B$ onto a normal $A$ with equal
positive-length vectors, there is $z\ne0$ with
\begin{align}
    VB^{\mathbf w}
        &=z\,V'B^{\mathbf w},
    &B^{\mathbf w}W
        &=z^{-1}B^{\mathbf w}W'
    \label{eq:fbc25_dressed_comparison}
\end{align}
for every word $\mathbf w$ longer than twice the nilpotency length.  Neither
source shows $V=zV'$ itself.  The scalars $\beta_{g,h}$ and $\gamma_{g,x}$
relating two choices of tensors, and the $L$-symbols, are therefore not
defined as bare scalars by the printed argument.

\section{The local fix}

All three families are defined by the dressed comparison
(\ref{eq:fbc25_dressed_comparison}): $\beta_{g,h}$ compares the left fusion
tensors of two choices against long words of $O_gO_h$, $\gamma_{g,x}$
compares the left action tensors against long words of $O_g\cdot A_x$, and
$L^x_{g,h}$ compares the two left boundaries against long words of
$(O_gO_h)\cdot A_x$.  Each is unique, because the left boundary of a
reduction onto a normal tensor of positive bond dimension does not vanish on
arbitrarily long words.

With these definitions the gauge law (\ref{eq:fbc25_dressed_l_gauge}) holds
exactly.  Changing the action tensors of $h$ and then of $g$ multiplies the
first boundary by $\gamma_{h,x}\gamma_{g,hx}$ against long words, after
pulling the second comparison back along the reduction by the action of $h$.
Changing the fusion tensor and then the action tensor of $gh$ multiplies the
second boundary by $\beta_{g,h}\gamma_{gh,x}$.  Comparing the two
expressions for the first boundary of the original choice on a word on which
it does not vanish gives (\ref{eq:fbc25_dressed_l_gauge}).

Conversely, rescaling actual tensors as in
(\ref{eq:fbc25_dressed_fusion_gauge}) and
(\ref{eq:fbc25_dressed_action_gauge}) is again a choice of fusion and action
tensors, and its dressed comparison scalars are exactly $\beta$ and
$\gamma$.  Hence every joint scalar gauge of the $L$-symbols is realized by
tensors.  This is what the later arguments of the source use: the single-block
statement that $\omega=L=1$ in some gauge (line~1926), the stabilizer form of
nonanomalous $L$-symbols (\texttt{lemma:h1h2}, lines~7028--7042), and the
block-independence criterion (\texttt{prop:BI\_psi}, lines~7042--7064).  All
three hold verbatim for the dressed scalars.

Recovering bare scalars, that is $F'^<=\beta^{-1}F^<$ for every pair of
choices, would require an argument absent from both sources and is not
asserted.\footnote{Formal declarations:
\leanid{MPOTensor.GroupFamily.FusionData.rescale},
\leanid{MPOTensor.GroupFamily.BlockActionData.relativeScalar},
\leanid{MPOTensor.GroupFamily.BlockActionData.lSymbol\_eq\_gauge}, and
\leanid{MPOTensor.GroupFamily.BlockActionData.lSymbol\_rescale} in
\doclink{TNLean/MPS/Symmetry/MPOSymmetry/PermutedBlocksGauge}.}

\textbf{Verdict: resolved local fix.}
The dressed scalars satisfy every printed relation, and every printed gauge is
realized by rescaled tensors.

\bibliographystyle{alpha}
\bibliography{references}

\end{document}
